% Generated draft: requires mathematical and source review.
\documentclass[11pt]{article}
\usepackage[margin=25mm]{geometry}
\usepackage{fontspec}
\setmainfont{Latin Modern Roman}
\usepackage{amsmath,amssymb,mathtools,unicode-math}
\setmathfont{Latin Modern Math}
\usepackage{enumitem,needspace,xurl,hyperref}
\setlength{\parindent}{0pt}
\setlength{\parskip}{5pt}
\setlength{\emergencystretch}{3em}
\newcommand{\N}{\mathbb{N}}
\newcommand{\Z}{\mathbb{Z}}
\newcommand{\Q}{\mathbb{Q}}
\newcommand{\R}{\mathbb{R}}
\newcommand{\C}{\mathbb{C}}
\newcommand{\F}{\mathbb{F}}
\newcommand{\A}{\mathbb{A}}
\newcommand{\PP}{\mathbb{P}}
\newcommand{\E}{\mathbb{E}}
\newcommand{\eps}{\varepsilon}
\newcommand{\dd}{\,\mathrm{d}}
\begin{document}
\section*{UnsolvedMath ID: 6000006}
\textbf{Source problem code:} catalogue entry
\subsection{Definitions and mathematical statement}
An affine flat manifold is a smooth manifold with a torsion-free flat affine connection, equivalently an affine atlas whose transition maps are restrictions of elements of $\operatorname{Aff}(\mathbb R^n)$. This does not mean a manifold with a flat Riemannian metric. Determine which groups occur as the fundamental group $\pi_1(M)$ of a compact affine flat manifold $M$.

The supplied summary singles out the complete case. Completeness here means geodesic completeness of the flat affine connection. A connected complete affine manifold can be presented as $\mathbb R^n/\Gamma$ for a free properly discontinuous affine action, and compactness is cocompactness. In that setting the cited Auslander conjecture asks whether $\Gamma$ is virtually polycyclic: it has a finite-index subgroup admitting a finite subnormal series with cyclic successive quotients. Thus its group-theoretic target is
\[
\Gamma\leq\operatorname{Aff}(\mathbb R^n),\quad
\Gamma\text{ acts freely, properly discontinuously and cocompactly}
\quad\Longrightarrow\quad\Gamma\text{ is virtually polycyclic}.
\]
The original question does not assume completeness, so it is broader than this explicitly complete formulation.

The supplied summary reports the complete-case conjecture proved for dimensions at most six by Abels--Margulis--Soifer and open in general dimension. Its parenthetical description of virtually polycyclic groups includes ``polynomial growth,'' which is not a consequence of virtual polycyclicity in general; that extra growth condition is not added to the conjecture. Its mention of noncompact Margulis spacetimes is background about a different quotient class.
\subsection{Short English statement}
Characterize fundamental groups of compact affine flat manifolds. Distinguish the unqualified question from Auslander's virtually-polycyclic conjecture for complete compact affine manifolds, and do not add an unsupported polynomial-growth requirement.
\subsection{Sources}
\begin{itemize}
\item Furuhata, Matsuzoe, and Urakawa, \emph{Open Problems in Affine Differential Geometry and Related Topics} (1998), item 2(a), scanned PDF pages 1--2, \url{https://www.jstage.jst.go.jp/article/iis/4/2/4_2_125/_pdf/-char/en}; accessed 2026-07-29, manual OCR from page images recorded in the provenance. The proposer field reads ``Yoe Itokawa; attributed to D. Gromoll.''
\item Auslander (1964), as the complete-case conjecture attribution in the supplied summary; no publication details are given. Abels--Margulis--Soifer, arXiv:1211.2525, is the supplied dimension-at-most-six reference. Margulis spacetimes are mentioned without a separate citation.
\item Local registry \texttt{data/unsolvedmath/problems.json}, permanent UnsolvedMath ID $6000006$, source code catalogue entry, \emph{Fundamental Groups of Compact Affine Flat Manifolds}. Dataset attribution: Ulam AI and the UnsolvedMath Contributors.
\item The supplied rights note says the J-STAGE article page marks this article CC BY 4.0.
\end{itemize}
\end{document}
